\documentclass{article}
\usepackage[utf8]{inputenc}
\usepackage[T1]{fontenc}
\usepackage{geometry}
\geometry{
    a4paper,
    left=2.5cm,
    right=2.5cm,
    top=2.5cm,
    bottom=2.5cm
}
\usepackage{graphicx}
\usepackage{float}
\usepackage{hyperref}
\usepackage[english]{babel}
\usepackage{amsmath}
\usepackage{amssymb}
\usepackage{array}
\usepackage{booktabs}
\usepackage{caption}
\usepackage{subcaption}
\usepackage{enumitem}
\usepackage{fancyhdr}
\usepackage{titlesec}
\usepackage{url}
\usepackage{textcomp}
\usepackage{xcolor}
\usepackage{listings}
\usepackage{multirow}
\usepackage{longtable}
\usepackage{tabularx}
\usepackage{times}

% Set up page style
\pagestyle{fancy}
\fancyhf{}
\rfoot{\thepage}
\renewcommand{\headrulewidth}{0pt}

% Title formatting
\titleformat{\chapter}[display]{\normalfont\bfseries}{\chaptertitle\thechapter}{0pt}{\Huge}
\titlespacing*{\chapter}{0pt}{50pt}{40pt}

% Document info
\title{Template Document}
\author{Academic Support Agent}
\date{\today}

\begin{document}
\maketitle

% cr ation de la page de titre
\maketitle

%%%%%%%%%%%%%%%%%%%%%%%%%%%%%%%%%%%%%%%%%%%%%%%%%%%
%% remerciements
%%%%%%%%%%%%%%%%%%%%%%%%%%%%%%%%%%%%%%%%%%%%%%%%%%%

\chapter*{Remerciements}
\thispagestyle{empty}
Je tiens   remercier Machin Bidule  pour ses conseils avis s.J'adresse  galement toute ma reconnaissance   Madame Truc Machin ainsi qu'  ses coll gues pour m'avoir accept  dans leur cercle tr s r serv . \\

Blabla

% cr ation de la table des mati res
\MyToc


\newpage        % saut de page

\newpage        % saut de page
% Justification moins stricte : des mots ne d passeront pas des paragraphes
\sloppy          

%%%%%%%%%%%%%%%%%%%%%%%%%%%%%%%%%%%%%%%%%%%%%%%%%%%
%% Introduction
%%%%%%%%%%%%%%%%%%%%%%%%%%%%%%%%%%%%%%%%%%%%%%%%%%%
% La commande '\section' marque le d but d'une nouvelle section.
% Elle prend le titre de la section en param tre.
% Latex se charge d'effectuer la num rotation automatiquement.

\section{Introduction}

L'utilisation de logiciels tels que MS Word ou LibreOffice pour la r daction de rapport aboutit g n ralement   des documents tr s h t rog nes,   la lisibilit  douteuse, ce qui entra ne des difficult s lors de leur  valuation. Pour cette raison, c'est \href{http://fr.wikipedia.org/wiki/LaTeX}{\LaTeX} \cite{LaTeXProject2010, Wikipedia2011a} qui a  t  retenu pour les rapports de cette UE.

% Pour cr er un nouveau paragraphe, il faut laisser une ligne vide apr s le paragraphe pr c dent.
La mise en forme d'un document est aussi importante que son contenu. Il est clair qu'un rapport bien pr sent  mais dont le contenu est inadapt  ou incorrect obtiendra une note faible. Mais de la m me fa on, un rapport dont le contenu est excellent et dont la pr sentation est mauvaise r coltera probablement une mauvaise note, car il sera difficile d'en comprendre les id es. Donc pour r sumer, une bonne pr sentation est une condition n cessaire, mais pas suffisante pour produire un bon rapport. Tout  a pour dire que la mise en forme de vos rapports sera  valu e au m me titre que leur contenu.
	
% La commande '\LaTeX' utilis e ci-dessous permet d' crire le mot "LaTeX" dans sa mise en forme officielle.
% La commande '\href{url}{texte}' utilis e ci-dessous permet d'inclure un hyperlien dans le document, en l'associant   un texte diff rent de l'adresse elle-m me.
% La commande '\url{url}' affiche  galement un hyperlien, mais cette fois il est associ    l'adresse elle-m me (et non pas   un autre texte).
% La commande '\texttt{texte}' utilis e ci-dessous permet d'ins rer du texte avec une police de caract res de type "listing" (i.e. courrier).
\LaTeX{}, pour simplifier, est un langage permettant de programmer la mise en forme d'un document. Donc, au lieu de composer votre document dans un traitement de texte, vous allez  crire seulement sa structure et son contenu sous la forme d'un code source (extension \texttt{.tex}). 
	
% Comme dans le paragraphe ci-dessous, vous pouvez utiliser la commande '\^a' 
% pour ins rer des accents circonflexes (ici sur la lettre 'a').
% De la m me fa on, la commande \`e permet d'ins rer un accent grave, comme dans ' ',
% tandis que \'e ins rera un accent aigu comme dans ' '.
\LaTeX{} et les outils associ s sont libres (et gratuits). Il s'agit d'un standard dans le domaine de la recherche acad mique : la plupart des articles soumis   des conf rences ou   des journaux scientifiques sont mis en forme gr\^ace   ce syst me.
	
% La commande '\cite{cl }' permet de citer une r f rence bibliographique. Pour en tenir compte,
% le document devra  tre compil  avec BibTeX, puis encore avec PDFLaTeX. A noter que les utilisateurs
% de Texmaker doivent modifier l'option associ e   BibTeX de mani re   obtenir "bibtex %"   la place de "bibtex %.aux"
Bien s\^ur, le but de ce cours n'est pas d'apprendre   utiliser \LaTeX{}, c'est pourquoi le pr sent document tient lieu   la fois de tutoriel et de mod le pour la r daction de M moires. Il ne s'agit pas d'un manuel sur \LaTeX{} : de nombreuses ressources disponibles sur le Web vous apprendront tout ce que vous devez savoir sur les bases de \LaTeX{}. Si vous n'avez jamais utilis  \LaTeX{}, il est recommand  de lire un texte d'introduction tel que ces WikiBooks en fran ais \cite{Wikibooks2010} ou en anglais \cite{Wikibooks2011}, en compl ment   ce document.
	
Le code source \LaTeX{} de ce document vous est fourni afin que vous l'utilisiez comme base pour vos propres rapports. Il est abondamment comment  afin d'en faciliter la compr hension. Il est recommand  de lire la version PDF et la version \LaTeX{} simultan ment pour bien en comprendre le fonctionnement. Pour  crire votre propre rapport, il est recommand  de faire une copie du fichier \texttt{.tex} et de s'en servir de base apr s avoir effac  les parties inutiles et les commentaires.
	  


%%%%%%%%%%%%%%%%%%%%%%%%%%%%%%%%%%%%%%%%%%%%%%%%%%%
%% Installation et configuration
%%%%%%%%%%%%%%%%%%%%%%%%%%%%%%%%%%%%%%%%%%%%%%%%%%%
\section{Mise en place du syst me}
On peut distinguer deux fa ons de produire un document \LaTeX{} : soit via une application locale, soit gr ce   une application Web. Les trois premi res sous-sections suivantes sont consacr es   la premi re approche, alors que la derni re porte sur l'utilisation de la seconde.

\subsection{Application Web}
Si vous ne voulez pas installer et configurer \LaTeX{} (ainsi que l' diteur qui va avec) sur votre syst me, vous pouvez vous tourner vers une application Web du type \href{https://www.overleaf.com/signup?ref=d62cb1694be6}{OverLeaf}. Vous devez d'abord vous cr er un compte (c'est gratuit), puis vous pourrez d finir et compiler vos propres documents en ligne. 

%Notez que la source \LaTeX{}  du pr sent document est disponible sur OverLeaf   l'adresse : \url{https://www.overleaf.com/latex/templates/modele-rapport-uapv/pdbgdpzsgwrt\#.Vo_iNvnhDmF}. Il s'agit d'un mod le   partir duquel vous pouvez instancier votre propre document.

L'inconv nient principal de l'application en ligne (par rapport   l'application locale) est que la compilation prend sensiblement plus de temps. Par contre, les avantages sont multiples :  pas d'installation requise, simplicit  d'utilisation, utilisable sur la plupart des navigateur, etc. De plus, OverLeaf permet le travail collaboratif : partage de documents, suivi des modifications, gestion des versions, syst me de commentaires, etc. Diff rents tutoriels sont disponibles sur le site officiel.


% La commande '\subsection' fait la m me chose que \section, mais pour le niveau inf rieur.
\subsection{Application "software"}
\subsubsection{Installation}
\label{sec:installation}
% La commande '\@.' permet de traiter un point qui ne termine pas une phrase 
% (utile pour les abr viations comme 'etc.').
Le syst me \LaTeX{}  tant libre, il est disponible sur les principaux syst mes d'exploitation (Unix, Linux, Windows, MacOS, etc\@.). Il est extr mement directement int gr  sous Linux par d faut.
% La command '\textit{texte}' utilis e ci-dessous permet de mettre du texte en italique.
La version de base d'Ubuntu contient le package \texttt{TexLive2009}, qui est une distribution \LaTeX{} suffisante pour g n rer vos rapports. Elle contient, entre autres, diff rentes extensions de \LaTeX{} utilis es dans ce document. V rifiez dans votre gestionnaire de packages que \texttt{TexLive} est bien install , et si ce n'est pas le cas : installez-le.

Si vous utilisez Windows, la distribution \LaTeX{} habituellement utilis e est \href{http://miktex.org/}{MiKTeX} \cite{MiKTeX2015}. Une alternative peut  tre d'utiliser une machine virtuelle telle que \href{http://www.virtualbox.org/}{VirtualBox} \cite{Oracle2011} pour installer Ubuntu sur votre ordinateur. La machine virtuelle va simuler le comportement d'un ordinateur, ce qui vous permettra d'ex cuter Ubuntu depuis un autre OS. L'inconv nient est que votre machine doit \^etre assez puissante pour pouvoir g rer les calculs suppl mentaires d\^us   l' mulation. Autrement dit :  a sera peut- tre un peu lent sur un PC portable. 

Quel que soit le syst me d'exploitation, vous avez besoin d'un package \LaTeX{} permettant d' crire des documents en fran ais. Il s'agit du package \texttt{texlive-lang-french}, que vous pouvez installer depuis le gestionnaire de packages lui aussi (Ubuntu) ou   partir du gestionnaire de packages de MiKTeX (Windows). Sans ce package, vous ne pourrez pas compiler le document.

Vous avez aussi besoin d'un  diteur de texte. Vous pouvez utiliser n'importe lequel, peu importe. Par exemple, \texttt{gedit} qui est g n ralement install  avec Ubuntu convient tout   fait. Cependant, pour plus de confort, il est tr s fortement recommand  d'installer plut t un  diteur adapt    \LaTeX{}. Le logiciel libre \href{http://www.xm1math.net/texmaker/}{Texmaker} \cite{Texmaker2011} est suffisant, et c'est lui qui sera utilis  dans les descriptions donn es dans ce document. A noter que ce logiciel existe aussi pour Windows et d'autres OS. Il apporte diff rentes fonctionalit s bien pratiques, comme la possibilit  de facilement compiler les documents. Sous Ubuntu, vous pouvez l'installer depuis le gestionnaire de packages. Notez que le logiciel que vous utiliserez pour  diter vos documents devra obligatoirement disposer d'un correcteur orthographique : c'est le cas de Texmaker (cf\@. la section \ref{sec:configuration} pour savoir comment le configurer).

Enfin, il est possible (et recommand ) d'utiliser un logiciel sp cifique pour g rer la bibliographie, tel que \href{http://jabref.sourceforge.net/}{JabRef} \cite{JabRef2008}. Ce logiciel libre peut  tre install    partir du gestionnaire de packages (il existe aussi pour Windows ou Mac OS, car il est programm  en Java). Il permet d'ajouter/supprimer/ diter des r f rences bibliographiques, et de g n rer un fichier au format BibTeX (voir la section \ref{sec:biblio}).
	
\subsubsection{Configuration}
\label{sec:configuration}
% La commande '\ieme' permet d'ins rer l'abr viation fran aise pour 'deuxi me', 'troisi me', etc. : '2 me', '3 me', etc.
% Le caract re d' chappement '\' permet de d sactiver les caract res sp ciaux tels que "%", comme ci-dessous.
% La commande \ref{nomlabel} permet de faire r f rence   un label cr  avec la commande '\label{nomlabel}' 
% (ici c'est une figure). Un hyperlien sera automatiquement plac  dans le document PDF r sultant de la compilation.
L'essentiel du travail de configuration pour \LaTeX{} et Texmaker a  t  effectu  lors de l'installation, donc il ne reste pas grand chose   faire. Vous devez d'abord modifier une option de Texmaker. Allez dans (\textit{Options / Configure Texmaker}) et localisez l'option de BibTeX (3\ieme ligne). Modifiez cette option de mani re   avoir \texttt{"bibtex \%"}   la place de \texttt{"bibtex \%.aux"} (voir la figure \ref{fig:bibtex}).

% La commande '\begin{figure}' permet d'ins rer une figure.
% Latex place automatiquement la figure dans le texte, en fonction des lettres pass es en param tres (ici : htb) :
% Le h signifie exactement que Latex a le droit de la placer exactement ici (h=here), le t signifie en haut d'une page (t=top) et le b en bas (b=bottom). Attention, s'agit seulement d'autorisations et non pas d'obligations.
% Si vous ne savez pas quoi mettre, laissez 'htb'.
\begin{figure}[htb]
	% La commande '\centering' va centrer la figure
	\centering
	% la commande \fbox permet de placer une bordure autour de l'image (optionnel).
	% La commande '\includegraphics' permet d'ins rer en tant que figure une image contenue dans un fichier 
	% (ici : fichier jpeg). Son option scale permet de r duire (<1) ou agrandir (>1) l'image.
	\fbox{\includegraphics[scale=0.4]{images/bibtex.jpg}}
	% La commande '\caption' permet de d finir la l gende de la figure (obligatoire).
	% L'option entre crochets ([texte]) est une version courte de la l gende, utilis e pour g n rer la liste des figures.
	\caption{Il est n cessaire de modifier l'option de BibTeX dans Texmaker.}
	% La commande '\label' donne un nom   la figure, et permet d'y faire r f rence dans le texte, 
	% grace   la commande '\ref', d j  utilis e pr c demment.
	\label{fig:bibtex}
\end{figure}

Vous devez ensuite configurer Texmaker pour qu'il utilise le dictionnaire fran ais lors de la v rification orthographique. Allez dans \textit{Options / Configure Texmaker / Editor / Spelling Dictionnary} et cliquez sur le bouton pour s lectionner le fichier \texttt{fr\_FR.dic} (c\@.f\@. figure \ref{fig:dictionnaire}). Si ce fichier n'est pas disponible, vous devez installer depuis le gestionnaire de packages ou MiKTeX un package dont le nom est du type \texttt{myspell-fr}. Une fois le dictionnaire s lectionn  dans Texmaker, les fautes d'orthographes seront soulign es en rouge. Attention, le logiciel ne d tecte pas les fautes de grammaire ou d'accord (genre, pluriel, conjugaison...). Soyez-donc attentifs   ce type d'erreurs !

\begin{figure}[htb]
	\centering
	\fbox{\includegraphics[scale=0.4]{images/dictionnaire.jpg}}
	\caption{Configuration du dictionnaire utilis  par le v rificateur d'orthographe.}
	\label{fig:dictionnaire}
\end{figure}

Pour pouvoir g n rer directement un fichier PDF, vous devez s lectionner le compilateur appropri . Dans la barre d'outil de Texmaker, cliquez sur \textit{LaTeX} et s lectionnez   la place \textit{PDFLaTeX}. Quand vous aurez besoin de compiler un document avec BibTeX (cf. section \ref{sec:biblio}), proc dez de la m me mani re avec l'option BibTeX.

\subsection{Compilation}
Texmaker facilite grandement la cr ation de documents \LaTeX{}. Il faut d'abord cr er un nouveau fichier, y  crire le code source \LaTeX{} et l'enregistrer. Un exemple de code minimal est donn  ci-dessous, que vous pouvez utiliser pour tester votre installation :
% Cette commande \vspace permet d'ins rer un espace. On l'utilise ici pour bien
% s parer le paragraphe pr c dent de l'exemple.
% N'utilisez pas cette commande dans vos rapports ! Laissez LaTeX se
% charger de faire la mise en page tout seul.
\vspace{0.3cm}\\
\noindent{}\texttt{\textbackslash{}documentclass\{article\}}\\
\noindent{}\texttt{\textbackslash{}begin\{document\}}\\
\indent{}\texttt{Ceci est un document.}\\
\noindent{}\texttt{\textbackslash{}end\{document\}}
\vspace{0.3cm}

Il suffit ensuite de cliquer sur la fl che situ e   gauche de \textit{PDFLaTeX} dans la barre d'outils de Texmaker, comme indiqu  dans la figure \ref{fig:compilation}, pour lancer la compilation du document. Vous pouvez  galement tester votre installation sur le fichier \texttt{.tex} fourni avec ce document : vous devez obtenir exactement le m me PDF, puisqu'il a  t  g n r  de cette fa on.

\begin{figure}[htb]
	\centering
	\fbox{\includegraphics[scale=0.4]{images/compilation.jpg}}
	\caption{Le fait de cliquer sur la fl che indiqu e d clenche la compilation du document \LaTeX{}.}
	\label{fig:compilation}
\end{figure}

Si la compilation rencontre des erreurs, celles-ci seront affich es dans la fen\^etre \textit{Messages / Log File} situ e en bas de Texmaker (figure \ref{fig:log}). Sinon, vous pourrez visualiser le document PDF en cliquant sur la fl che situ e   gauche de \textit{View PDF} (dans la barre d'outils). Si, lors lors du premier test, le compilateur ne trouve pas le fichier de classe \texttt{ceri.cls}, v rifiez que vous l'avez bien plac  dans le m me dossier que votre fichier \texttt{.tex}.

\begin{figure}[htb]
	\centering
	\fbox{\includegraphics[scale=0.4]{images/log.jpg}}
	\caption{Texmaker affiche le r sultat de la compilation dans la fen\^etre situ e sous l' diteur de texte. En cas d'erreur, les fl ches bleues situ es   gauche permettent de visualiser les messages.}
	\label{fig:log}
\end{figure}

A noter qu'il est possible que le compilateur indique   l'aide d'un avertissement (warning) la n cessit  de compiler une autre fois le document. C'est notamment le cas si des r f rences bibliographiques ont  t  modifi es (voir section \ref{sec:biblio}) ou si la structure du document a chang  (cf\@. section \ref{sec:structure}).



%%%%%%%%%%%%%%%%%%%%%%%%%%%%%%%%%%%%%%%%%%%%%%%%%%%
%% G n ralit s
%%%%%%%%%%%%%%%%%%%%%%%%%%%%%%%%%%%%%%%%%%%%%%%%%%%
\section{G n ralit s}
Vous devez indiquer le titre du document gr ce   la commande \texttt{\textbackslash{}title} d finie en d but de fichier. Les noms des auteurs sont   sp cifier dans la commande \texttt{\textbackslash{}author}.
		
\subsection{Mise en forme de base}
\label{sec:miseenforme}
La mise en forme de base est d crite clairement dans les r f rences d j  cit es \cite{Wikibooks2010, Wikibooks2011}. Voici rapidement quelques  l ments utiles (r f rez-vous au code source pour comprendre le fonctionnement de ces commandes).
	
Le texte peut  tre mis en \textbf{gras} avec \texttt{\textbackslash{}textbf} ou en \textit{italique} avec \texttt{\textbackslash{}textit}. Il est possible d'inclure des listes non-num rot es avec \texttt{\textbackslash{}begin\{itemize\}} et \texttt{\textbackslash{}item}, telles que :	
% La commande '\begin{itemize}' permet de d finir une liste dont les  l ments ne sont pas num rot s.
\begin{itemize}
	% chaque  l ment est d fini avec la commande '\item'
	\item premier  l ment ;
	\item 2\ieme  l ment ;
	\item troisi me  l ment.
	\begin{itemize}
		\item on peut aussi placer des sous-listes ;
		\begin{itemize}
			\item et rajouter encore d'autres niveaux ;
			\item etc.
		\end{itemize}
		\item on continue au niveau sup rieur ;
	\end{itemize}
	\item et au niveau encore sup rieur.
\end{itemize}

Ou bien des listes num rot es avec \texttt{\textbackslash{}begin\{enumerate\}}, comme :
% La commande '\begin{enumerate}' fonctionne de la meme fa on qu'itemize, mais les  l ments sont num rot s.
\begin{enumerate}
	% Comme pour itemize, chaque  l ment est d fini avec la commande '\item'
	\item premier  l ment ;
	\item 2\ieme  l ment ;
	\item troisi me  l ment.
	\begin{enumerate}
		\item l  encore, on peut placer des sous-listes ;
		\item  a marche exactement pareil que pour les listes non-num rot es ;
		\begin{itemize}
			\item on peut m me m langer les deux types de listes comme ici ;
			\item et ici.
		\end{itemize}
	\end{enumerate}
\end{enumerate}
	
% La commande '\colorbox{couleur}' permet de modifier la couleur de fond du texte.
Il est aussi possible de placer un fond color  derri re le texte avec la commande \texttt{\textbackslash{}colorbox\{couleur\}}. Pour les couleurs, vous pouvez soit utiliser le nom de la couleur en anglais, par exemple \colorbox{red}{\texttt{red}}, soit d finir une couleur pr cise en utilisant des valeurs num riques.

\textbf{Remarque :}  vitez au maximum les notes de bas de page (pas de commande \texttt{\textbackslash{}footnote}) : placez vos remarques directement dans le texte, entre parenth ses.

\subsection{Structure}
\label{sec:structure}
Pour cr er des paragraphes s par s, vous devez sauter une ligne (voir le code source). La structure du document est d finie en utilisant des commandes sp cifiques. Par ordre de niveau d croissant, on a :
\begin{enumerate}
	\item \texttt{\textbackslash{}section\{Titre de la section de niveau 1\}}
	\item \texttt{\textbackslash{}subsection\{Titre de la section de niveau 2\}}
	\item \texttt{\textbackslash{}subsubsection\{Titre de la section de niveau 3\}}
	\item \texttt{\textbackslash{}paragraph\{Titre de la section de niveau 4\}}
	\item \texttt{\textbackslash{}subparagraph\{Titre de la section de niveau 5\}}
\end{enumerate}

Les deux derniers niveaux ne sont pas utilis s dans ce document, car ils n'y sont pas n cessaires. Pour l'exemple, les voici :
\subsubsection{Section de niveau 3}
(Il est n cessaire de mettre une section de niveau 3 entre la section de niveau 2 et celle de niveau 4, sinon une section factice de niveau 3 num rot e 0 sera automatiquement cr e).
\paragraph{Section de niveau 4}
Le niveau 4 n'est pas vraiment une section mais plut t un titre de paragraphe. Pour cette raison, le titre n'est pas sur une ligne s par e, mais il est directement plac  dans le paragraphe concern .
\subparagraph{Section de niveau 5}
M me chose pour le niveau 5 : on a un titre ins r  en d but de paragraphe.
	
\textbf{Remarques :} ne d finissez pas une section seulement pour y placer 1 ou 2 lignes de texte. Une section doit avoir un contenu assez long, sinon ce n'est pas la peine d'en cr er une.
	
\subsection{R f rences crois es}
\label{sec:references}
\LaTeX{} permet de d finir un label pour marquer un endroit particulier du document. Il est ensuite possible de faire r f rence   ce label. La r f rence appara tra dans le texte sous la forme d'un lien hypertexte, ce qui est pratique pour naviguer dans le document. 
	
Un label se d finit avec la commande \texttt{\textbackslash{}label\{nomlabel\}} et une r f rence   un label avec \texttt{\textbackslash{}ref\{nomlabel\}}. Si le label correspond   un objet num rot  (figure, section, table...), alors c'est son num ro qui sera affich .
	
La convention veut qu'on utilise des pr fixes sp cifiques dans les noms des labels en fonction du type d'objet r f renc . Vous devez utiliser cette convention dans vos rapports (voir la table \ref{tab:prefixes}).
	
% La commande '\begin{table}' permet d'ins rer une table dans le document.
% Elle fonctionne plus ou moins comme '\figure', notamment en ce qui concerne
% son positionnement avec h, t et b.
\begin{table}[htb]
	% La commande '\centering' sert encore   centrer la table.
	\centering
	% Cette commande permet de colorer une ligne sur deux, pour am liorer la visibilit 
	\rowcolors{1}{LightColor}{}
	% La commande '\tabular' cr  la table proprement dite.
	% le param tre liste les colonnes de la table. Chaque lettre repr sente une colonne
	% et son alignement : l=align e   gauche, r=align e   droite et c=centr e.
	% Ici, on a deux colonnes align es   gauche.
	\begin{tabular}{l l}
		\hline
		% La commande '\rowcolor' colorie la ligne de titre de couleur plus fonc e.
		\rowcolor{DarkColor} 
		% Le contenu de la table est d fini ligne   ligne. Les colonnes
		% sont s par es par le caract re '&' et la fin de la ligne est 
		% indiqu e par la chaine '\\' (obligatoire).
		% La ligne de titre doit etre en gras (commande '\textbf')
		\textbf{Type d'objet}	& \textbf{Pr fixe}	\\
		% La commande '\hline' ins re une s paration d' paisseur normale.
		\hline
		section					& \texttt{sec}	\\
		figure 					& \texttt{fig}	\\
		table 					& \texttt{tab}	\\
		equation 				& \texttt{eq}	\\
		\hline
	\end{tabular}
	% La l gende est d finie comme pour une figure.
	\caption{Pr fixes utilis s conventionnellement pour d finir les noms des labels.}
	% Le label aussi est d fini comme pour une figure, mais en utilisant le pr fixe 'tab' plutot que 'fig'.
	\label{tab:prefixes}
\end{table}
	
Il est absolument obligatoire que tout objet ( quation, table, figure...) num rot  soit cit  dans votre texte. Donc chacun doit avoir son label, et ce label doit apparaitre au moins une fois dans le texte. Bien s r vous ne devez pas vous contenter de citer chaque objet : ils doivent  tre suffisamment d crits, et de fa on pr cise. Un objet qui n'est pas mentionn  dans le texte ne sert   rien.



%%%%%%%%%%%%%%%%%%%%%%%%%%%%%%%%%%%%%%%%%%%%%%%%%%%
%% Objets
%%%%%%%%%%%%%%%%%%%%%%%%%%%%%%%%%%%%%%%%%%%%%%%%%%%
\section{Objets}
\label{sec:objets}
\LaTeX{} permet d'ins rer un certain nombre d'objets num rot s : tables,  quations, figures et algorithmes. Ces objets doivent toujours poss der une l gende (sauf les  quations) et  tre num rot s. De plus, ils doivent  tre mentionn s, d crits et comment s dans le texte. Toute r f rence   un objet doit  tre faite en utilisant la commande \texttt{\textbackslash{}ref\{nomlabel\}} (cf\@. section \ref{sec:references}).
	 
\LaTeX{} se charge de placer les objets dans le document. Il est possible de restreindre le positionnement en pr cisant si on veut que l'objet en question soit plac    l'endroit o  il se trouve dans le document original (\texttt{h} pour \textit{here}), en haut de page (\texttt{t} pour \textit{top}) ou en bas de page (\texttt{b} pour \textit{bottom}). Il est possible de combiner ces valeurs si on veut donner plus de libert  au compilateur (ce qui est recommand ). Il est aussi possible que le compilateur place les objets ailleurs s'il est impossible de les placer   l'endroit sp cifi .
	
En plus du positionnement, les l gendes sont  galement g r es de fa on automatique. Elles sont d clar es dans les objets,   l'aide de la commande \texttt{\textbackslash{}caption\{l gende\}}.
	
\subsection{Tables}
\label{sec:tables}
La syntaxe permettant de d finir des tables en \LaTeX{} est relativement compliqu e. Vous vous limiterez   des tables de la forme d crite dans cette sous-section : cette mise en forme doit  tre respect e, afin de garder des tables lisibles. La commande \texttt{\textbackslash{}begin\{table\}[position]} permet d'ins rer une table dans le document. Le param tre \texttt{position} correspond   une combinaison de \texttt{h}, \texttt{t} et \texttt{b} qui permettra au compilateur de placer la table dans le document. Comme indiqu  auparavant, le mieux est d'utiliser les trois, c'est   dire \texttt{htb}.
	
Chaque table doit  tre centr e, ce qui est r alis  avec la commande \texttt{\textbackslash{}centering}. La premi re ligne de la table contient les titres des colonnes. Le texte contenu dans cette ligne de titre doit  tre en \textbf{gras}. Les traits de s paration sont dessin es avec \texttt{\textbackslash{}hline}.
	
Le contenu proprement dit de la table est d fini gr ce   la commande \texttt{\textbackslash{}begin\{tabular\}\{colonnes\}}. Le param tre \texttt{colonnes} permet de pr ciser   la fois le nombre de colonnes contenues dans la table, et leur alignement : 
\begin{itemize}
	\item \texttt{l} (\textit{left}) pour une colonne align e   gauche ;
	\item \texttt{r} (\textit{right}) pour qu'elle soit align e   droite ;
	\item \texttt{c} (\textit{center}) pour qu'elle soit centr e.
\end{itemize}

Les donn es de la table sont ensuite d crites ligne par ligne : la valeur de chaque colonne est s par e par un caract re \texttt{\&} et la fin de la ligne est marqu e par un \texttt{\textbackslash{}\textbackslash{}} (obligatoire).
	
Vos tables doivent ressembler exactement   celles donn es en exemple dans ce document : couleurs, s parations, l gende situ e en-dessous, etc. Les valeurs num riques doivent  tre align es   droite. Les r els doivent  tre d finis de mani re   avoir le m me nombre de d cimales. Les valeurs textuelles peuvent  tre align es   droite ou   gauche, ou bien centr es. Il est possible d'ins rer des  quations en-ligne (i\@.e\@. non-num rot es) dans une table (voir la section \ref{sec:equations}). La table \ref{tab:exemple} constitue un autre exemple   prendre en compte.
	
\begin{table}[htb]
	\centering
	\rowcolors{1}{LightColor}{}
	\begin{tabular}{l c r r c}
		\hline
		\rowcolor{DarkColor} 
		\textbf{Texte} 			& \textbf{Texte}			& \textbf{Entiers}	& \textbf{R els}	& \textbf{Formules}	\\ 
		\hline
		% Le caract re '$' est un d limiteur permettant de passer dans le mode math matique.
		% Le texte situ  entre les '$' est interpr t  comme une formule math matique et
		% est mis en forme de fa on sp cifique.
		Blabla bla bla bla bla	& Blabla bbla bla			& 21 				& 12,62				& $a+b=c$								\\
		Blabla bla bla bla		& Blabla bla bla bla bla	& 12356 			& 3456,21			& $x\cos{x}$							\\
		Blabla bla bla			& Blabla bla bbla bla		& 1 				& 45,87				& $\frac{\alpha^{2}+\beta}{\gamma}$	\\
		Blabla blabla bla bla	& Blabla bla bla bla		& 456 				& 45,16				& $\sqrt{x+y}$	\\
		Blabla bla blaa bla		& Blabla bla bla bla		& 78 				& 89,28				& $x_{1}-x_{2}$	\\
		Blabla bla bla bla bla	& Blabla bla bla bla		& 416 				& 10,10				& $\max_{x}{x}$	\\
		Blabla b bla bla		& Blabla blbla bla bla		& 9635 				& 99,99				& $\frac{x}{y}$	\\
		Blabla bla bbla bla		& Blabla bla b				& 965801 			& 46465,45			& $\Sigma_{x=1}^{10}{xy}$	\\
		Blabla bla				& Blabla bla				& 7 				& 48,78				& $\varphi^2$	\\
		\hline
	\end{tabular}
	\caption{Exemple de table contenant du texte, des valeurs enti res, r elles et des formules math matiques.}
	\label{tab:exemple}
\end{table}

Si votre table est haute et  troite, alors vous devez la d composer afin de mieux utiliser l'espace. Ceci est illustr  avec la table \ref{tab:etroit}, qui contient $50$ lignes pour seulement $2$ colonnes. Ceci est r alis  en utilisant l'environnement \texttt{minipage}, cf. le code source de ce document.
	
\begin{table}[htb]
	\rowcolors{1}{LightColor}{}
	% Chaque partie de la table doit  tre plac e dans un environnement "minipage".
	% Ici, le param tre "0.19\linewidth" indique que la premi re partie de la table
	% occupe 19% de la largeur de la page. Il faut donc adapter cette valeur en
	% fonction du nombre de parties   afficher.
	\begin{minipage}[t]{0.19\linewidth}
		\centering
		\begin{tabular}{l l}
			\hline
			\rowcolor{DarkColor} 
			\textbf{Col1} 			& \textbf{Col2} \\ 
			\hline
			1 & hdc \\
			2 & orf \\
			3 & seh \\
			4 & lig \\
			5 & jhg \\
			6 & azr \\
			7 & opi \\
			8 & hgj \\
			9 & xvc \\
			10 & ree \\
			\hline
		\end{tabular}
	\end{minipage}
	% Deuxi me partie de la table...
	\begin{minipage}[t]{0.19\linewidth}
		\centering
		\begin{tabular}{l l}
			\hline
			\rowcolor{DarkColor} 
			\textbf{Col1} 			& \textbf{Col2} \\ 
			\hline
			11 & ytr \\
			12 & xcv \\
			13 & fsd \\
			14 & pui \\
			15 & gsf \\
			16 & aze \\
			17 & bcv \\
			18 & wxc \\
			19 & uyr \\
			20 & ndz \\
			\hline
		\end{tabular}
	\end{minipage}
	% Troisi me partie de la table...
	\begin{minipage}[t]{0.19\linewidth}
		\centering
		\begin{tabular}{l l}
			\hline
			\rowcolor{DarkColor} 
			\textbf{Col1} 			& \textbf{Col2} \\ 
			\hline
			21 & ioe \\
			22 & snu \\
			23 & dgp \\
			24 & zbu \\
			25 & zvz \\
			26 & orv \\
			27 & czr \\
			28 & azh \\
			29 & cyh \\
			30 & ytj \\
			\hline
		\end{tabular}
	\end{minipage}
	% Quatri me partie de la table...
	\begin{minipage}[t]{0.19\linewidth}
		\centering
		\begin{tabular}{l l}
			\hline
			\rowcolor{DarkColor} 
			\textbf{Col1} 			& \textbf{Col2} \\ 
			\hline
			31 & zyc \\
			32 & opv \\
			33 & azv \\
			34 & xdf \\
			35 & yzd \\
			36 & sgg \\
			37 & sss \\
			38 & zdv \\
			39 & sgu \\
			40 & vzd \\
			\hline
		\end{tabular}
	\end{minipage}
	% Cinqui me partie de la table...
	\begin{minipage}[t]{0.19\linewidth}
		\centering
		\begin{tabular}{l l}
			\hline
			\rowcolor{DarkColor} 
			\textbf{Col1} 			& \textbf{Col2} \\ 
			\hline
			41 & oze \\
			42 & anb \\
			43 & kir \\
			44 & sfk \\
			45 & bed \\
			46 & pyj \\
			47 & zdj \\
			48 & pyj \\
			49 & zen \\
			50 & oef \\
			\hline
		\end{tabular}
	\end{minipage}
	\caption{Exemple de table haute et  troite, d compos e en 5 parties pour occuper plus efficacement l'espace disponible.}
	\label{tab:etroit}
\end{table}

\subsection{ quations}
\label{sec:equations}
Les  quations peuvent  tre introduites sous deux formes : \textit{en}-ligne et \textit{hors}-ligne. Une  quation en-ligne est int gr e au texte, par exemple $y=ax+b$. La syntaxe \LaTeX{} consiste simplement   l'entourer avec des caract res \texttt{\$} : le 1\ier~\texttt{\$} active ce qu'on appelle le \textit{mode math matique} et le 2\ieme~le d sactive. Le texte situ  entre les deux est interpr t  de fa on diff rente du texte normal. Voyez \cite{Wikibooks2010, Wikibooks2011} pour plus de d tails sur les nombreuses possibilit s de \LaTeX{} en ce qui concerne les formules math matiques, et le code source de la table \ref{tab:exemple} pour quelques exemples. Cette m thode permet aussi d'ins rer des formules dans les tables, comme par exemple dans la table \ref{tab:exemple}.

% La commande '\begin{equation}' permet de d finir une  quation hors-ligne.
% On ne donne pas de l gende aux  quations, et il est inutile de les centrer
% (elles le sont par d faut). On ne peut pas les positionner non plus (elles
% apparaissent l  o  elles sont d finies).
\begin{equation}
	f(s)=h(s)+g(s)
   	% Par contre, si on veut faire r f rence   l' quation, il est n cessaire
   	% d'y d finir un label.
	\label{eq:exemple}
\end{equation}

Les  quations hors-ligne sont s par es du texte et num rot es. En y associant un label, il est donc possible d'y faire r f rence dans le texte. Par exemple, l' quation ci-dessus est l' q. \ref{eq:exemple}. On les d finit en utilisant la commande \texttt{\textbackslash{}begin\{equation\}}. Une  quation n'a pas de l gende, et elle doit  tre d crite dans le texte qui y fait r f rence.

\textbf{Remarques :}
\begin{itemize}
	\item Vos variables et fonctions doivent avoir des noms courts : une lettre et un indice, maximum. Pas de nom du type ${mafonction}$, mais plut t $f$, $f'$ $f_1$, etc\@.
	\item Ne m langez pas les notations informatiques et math matiques. Par exemple n'utilisez pas la notation tableau (array), comme par exemple \texttt{tab[12]}, dans une  quation.
\end{itemize}
	
	
\subsection{Figures}
\label{sec:figures}
Une figure s'ins re de mani re tr s similaire   une table, except  qu'on utilise la commande \texttt{\textbackslash{}begin\{figure\}} au lieu de \texttt{\textbackslash{}begin\{table\}}. A part  a, la commande \texttt{\textbackslash{}centering} doit toujours  tre pr sente pour centrer la figure, la l gende doit  tre d finie en-dessous de la figure avec \texttt{\textbackslash{}caption}, et un label peut  tre cr  avec \texttt{\textbackslash{}label} pour r f rencer la figure, comme par exemple la figure \ref{fig:bibtex}.

\begin{figure}[th]
	\centering
	\includegraphics[scale=0.3]{images/AGROSUP_pre_sentation.JPG}
	\caption{AgroSup Dijon.}
	\label{fig:uapv}
\end{figure}

L'utilisation du compilateur \textit{PDFLaTeX} permet de traiter seulement des fichiers aux formats PNG, JPEG, GIF et PDF. L'image est sp cifi e au moyen de la commande \texttt{\textbackslash{}includegraphics[scale=s]\{fichier\}}, o  \texttt{fichier} est le chemin du fichier contenant l'image (les chemins relatifs sont exprim s par rapport   l'emplacement du fichier \texttt{.tex} en cours de compilation) et \texttt{s} est une valeur r elle indiquant l' chelle de l'image (optionnel). A la place du param tre \texttt{scale}, vous pouvez sp cifier la largeur ou la hauteur de l'image, par exemple avec \texttt{width=12cm} ou \texttt{height=5cm}.

Tous les algorithmes doivent  tre donn s sous forme de diagrammes du genre \href{http://en.wikipedia.org/wiki/Flowchart}{flowchart} \cite{Wikipedia2011}, qui appara tront dans le rapport en tant que figures. Il faut  viter   tout prix le pseudo-code. 

% \boitemagique{Titre de la boite}{
% Praesent placerat, ante at venenatis pretium, diam turpis faucibus arcu, nec vehicula quam lorem ut leo. Sed facilisis, augue in pharetra dapibus, ligula justo accumsan massa, eu suscipit felis ipsum eget enim.
% }

%%%%%%%%%%%%%%%%%%%%%%%%%%%%%%%%%%%%%%%%%%%%%%%%%%%
%% R f rences
%%%%%%%%%%%%%%%%%%%%%%%%%%%%%%%%%%%%%%%%%%%%%%%%%%%
\section{R f rences bibliographiques}
\label{sec:biblio}
Le syst me \LaTeX{} peut  tre coupl  avec BibTeX afin de g rer automatiquement les r f rences bibliographiques du document. La proc dure consiste d'abord   d finir un fichier de type BibTeX d crivant chaque r f rence en d tails. Comme indiqu  pr c demment, il est recommand  d'utiliser un logiciel sp cifique comme \href{http://jabref.sourceforge.net/}{JabRef} \cite{JabRef2008} pour r aliser cette t che. Le fichier \texttt{bibliographie.bib} fourni avec ce document est un exemple de fichier BibTeX. Il est plac  dans le m me dossier que le fichier \texttt{.tex}   compiler.

Une r f rence bibliographique doit toujours indiquer au moins les auteurs, le titre du document et l'ann e de publication. Dans un contexte acad mique, une r f rence vers un article de journal scientifique, comme \cite{Fortunato2010} par exemple, doit  galement contenir le nom du journal, le volume, le num ro et les pages. Pour une conf rence, comme par exemple \cite{Wei1989}, les champs additionnels sont le nom de la conf rence et les pages. Pour un rapport technique comme \cite{Rosvall2009a} ou une th se comme \cite{Gerbaud2010}, il faut aussi indiquer l'institution d'affiliation des auteurs.

Dans le fichier \texttt{.tex}, une r f rence est introduite au moyen de la commande \texttt{\textbackslash{}cite\{cle\}}, o  \texttt{cle} est la cl  associ e dans le fichier BibTeX   la r f rence d sir e. Par exemple la commande \texttt{\textbackslash{}cite\{Wikibooks2010\}} permet d'introduire la citation suivante : \cite{Wikibooks2010}. Il faut ensuite compiler le document avec BibTeX (voir la section \ref{sec:installation}), puis le recompiler avec PDFLaTeX.
	
Attention, lorsque vous d finissez une cl  n'utilisez pas de caract res accentu s. Vous pouvez mettre plusieurs r f rences dans une m me commande \texttt{\textbackslash{}cite}, par exemple \texttt{\textbackslash{}cite\{Wikibooks2010, Wikibooks2011\}} donne \cite{Wikibooks2010, Wikibooks2011}.
	
BibTeX et \LaTeX{} vont s'occuper de num roter les r f rences automatiquement. La commande \texttt{\textbackslash{}bibliography\{fichier\}} permet d'ins rer une liste des r f rences, g n ralement on la place   la fin du document. Le param tre \texttt{fichier} doit indiquer l'emplacement et le nom du fichier BibTeX   utiliser. Pour vous aider, le fichier \texttt{bibliographie.bib} est fourni   titre d'exemple.
	
Toutes les informations que vous int grez   votre rapport et qui n'ont pas  t  produites par vous-m me doivent  tre explicitement r f renc es. Vous devez donc indiquer tout document que vous utilisez, y compris ce qui est r cup r  sur internet. Cela vaut pour tout type d'information : document texte, images, vid os, code source, algorithmes, id es diverses.

Autrement dit, vous avez le droit de r utiliser absolument tout ce que vous d sirez. Par contre, vous devez   chaque fois indiquer exactement ce que vous r utilisez, et   qui/o  vous l'avez pris. De plus, vous devez expliquer ce que vous avez modifi  (par exemple s'il s'agit d'un algorithme) et la nature de ces modifications.



  
%%%%%%%%%%%%%%%%%%%%%%%%%%%%%%%%%%%%%%%%%%%%%%%%%%%
%% Conclusion
%%%%%%%%%%%%%%%%%%%%%%%%%%%%%%%%%%%%%%%%%%%%%%%%%%%
\section{Conclusion}
Avec \LaTeX{} on obtient un rendu impeccable mais il faut s'investir un minimum pour le prendre en main. Pensez   consulter le code source de ce document, qui est abondamment comment .
 
 
  
  
%%%%%%%%%%%%%%%%%%%%%%%%%%%%%%%%%%%%%%%%%%%%%%%%%%%
%% Bibliographie
%%%%%%%%%%%%%%%%%%%%%%%%%%%%%%%%%%%%%%%%%%%%%%%%%%%
\MyBibliography{bibliographie}

\begin{appendix}
    \chapter{Ici les Annexes A}
    \label{Annexes_A}

    \begin{center}
\begin{figure}[th]
	\centering
	\includegraphics[scale=0.3]{images/AGROSUP_pre_sentation.JPG}
	\caption{AgroSup Dijon.}
	\label{fig:uapv annexe A}
\end{figure}
    \end{center}

    \chapter{Seconde Annexe}
    \label{a_seconde}

    \begin{center}
       % \includegraphics[scale=0.28, angle=90]{./img/encartage.jpg}
    \end{center}

    \chapter{Troisi me annexe}
    \label{a_trois}

    \begin{center}
      %  \includegraphics[scale=0.28, angle=90]{./img/perso.jpg}
    \end{center}

    \chapter{Encore une annexe}
    \label{a_quatre}

    \begin{center}
       % \includegraphics[scale=0.5, angle=0]{./img/stage_prod.png}
    \end{center}
\end{appendix}
\end{document}